\section{Conclusion}
\label{sec:conclusion}

We have introduced Magic Gems, a geometric framework for representing magic squares 
as three-dimensional polyhedra.
This framework reveals deep connections between the classical algebraic definition of magic squares
and concepts from statistics, optimization, and statistical mechanics,
opening new avenues for understanding these ancient mathematical objects.

Our central theoretical contribution is the Complete Energy Characterization Theorem 
(\Cref{thm:complete_energy_char}), which establishes that an arrangement of integers in a square grid 
is a magic square if and only if its covariance with every row, column, and diagonal indicator vanishes.
Equivalently, magic squares are precisely the ground states of the complete energy functional 
$E_{\mathrm{full}}$, valid for all orders $n \geq 3$.
This characterization transforms the discrete, combinatorial magic property
into a statistical orthogonality condition with physical meaning.

The proof is elementary once the right framework is established:
each row/column covariance vanishes if and only if that line sums to $M(n)$,
and similarly for the diagonals.
The perturbation analysis confirms that magic squares are isolated local minima,
with every single-swap perturbation strictly increasing the energy.

We also investigated a natural simplification---the four-term ``low-mode'' energy 
using aggregate position coordinates $X$, $Y$ and diagonal indicators.
This characterizes magic squares for $n = 3$ (verified exhaustively),
but defines a strictly larger class of ``low-mode balanced'' arrangements for $n \geq 4$.
The existence of non-magic zeros of the low-mode energy, while initially surprising,
follows from a dimensional argument: four linear constraints cannot determine 
all $2n$ line-sum conditions.
Explicit counterexamples (\Cref{rem:four_term_insufficient}) and comparative energy analysis (\Cref{fig:energy_comparison})
illuminate this distinction,
highlighting the necessity of the complete row/column indicator basis for a rigorous characterization.

The energy landscape perspective offers the most compelling conceptual insight.
Magic squares are not merely configurations satisfying arithmetic constraints;
they are equilibrium states where an imbalance energy vanishes.
Like stones that have rolled to rest at the bottoms of valleys,
magic squares occupy the deepest wells in a rugged terrain of arrangements.
This view connects magic squares to optimization theory, statistical mechanics, and dynamical systems,
suggesting new algorithmic and theoretical approaches.

The Magic Gem framework opens several directions for future research.
The geometric classification of magic squares, particularly for $n = 4$ where 880 essentially different squares exist,
could complement traditional algebraic approaches.
Extensions to magic hypercubes, semi-magic squares, and pandiagonal squares 
would test the generality of our framework.
The ``low-mode balanced'' class identified by $E_{\mathrm{low}} = 0$ for $n \geq 4$
is itself an interesting mathematical object worthy of further study.

The Lo Shu square has fascinated humanity for nearly five thousand years,
its balanced arrangement of numbers inspiring mythology, art, and mathematics across cultures.
By transforming this arrangement into a three-dimensional object---a gem 
that sits at the bottom of an energy well, whose very shape encodes perfect equilibrium---we 
hope to have added a new facet to its enduring mathematical beauty.
